
\documentclass{article}
\usepackage{colortbl}
\usepackage{makecell}
\usepackage{multirow}
\usepackage{supertabular}

\begin{document}

\newcounter{utterance}

\centering \large Interaction Transcript for game `dond', experiment `semi\_en', episode 16 with Bentel rockers\_\_qwen3.5{-}4b\_\_Bentel\_iter\_3.
\vspace{24pt}

{ \footnotesize  \setcounter{utterance}{1}
\setlength{\tabcolsep}{0pt}
\begin{supertabular}{c@{$\;$}|p{.15\linewidth}@{}p{.15\linewidth}p{.15\linewidth}p{.15\linewidth}p{.15\linewidth}p{.15\linewidth}}
    \# & $\;$A & \multicolumn{4}{c}{Game Master} & $\;\:$B\\
    \hline

    \theutterance \stepcounter{utterance}  
    & & \multicolumn{4}{p{0.6\linewidth}}{
        \cellcolor[rgb]{0.9,0.9,0.9}{
            \makecell[{{p{\linewidth}}}]{
                \texttt{\tiny{[P1$\langle$GM]}}
                \texttt{You are playing a negotiation game in which you have to agree on how to divide a set of items among you and another player.} \\
\\ 
\texttt{Rules:} \\
\texttt{(a) You and the other player are given a set of items. Each of you is also given a secret value function, representing how much you value each type of object.} \\
\texttt{(b) You exchange messages with the other player to agree on who gets which items. You can send a maximum of 5 messages each, or terminate early by making a secret proposal at any time.} \\
\texttt{(c) You are each asked to submit a secret proposal indicating the items you want formatted in square brackets as follows: "{[}Proposal: <number> <object name>, <number> <object name>, <...>{]}"} \\
\texttt{(d) If your proposals are complementary, i.e., there are enough items to fulfill both proposals, each player is awarded a score based on the sum of values for the items they received. Otherwise, both of you get zero points.} \\
\texttt{(e) Your goal is to maximize the score you receive.} \\
\\ 
\texttt{Let us start.} \\
\\ 
\texttt{The set of available items is:} \\
\\ 
\texttt{1 bottle, 3 carrots, 2 glues} \\
\\ 
\texttt{Your secret value function is:} \\
\\ 
\texttt{bottle: 4, carrot: 0, glue: 3} \\
\\ 
\texttt{IMPORTANT: Your messages, unless it is the secret proposal, are directly transmitted to the other player, so do not include any response to the rules or text announcing your message. To make a secret proposal, use the indicated format. Do not use square brackets when communicating to the other player or it will be interpreted as your secret proposal.} \\
            }
        }
    }
    & \\ \\

    \theutterance \stepcounter{utterance}  
    & \multicolumn{4}{p{0.6\linewidth}}{
        \cellcolor[rgb]{0.9,0.9,0.9}{
            \makecell[{{p{\linewidth}}}]{
                \texttt{\tiny{[P1$\rangle$GM]}}
                \texttt{I value the bottle (4) and a glue (3) for a total of 7, while the carrots (0) add nothing to my score. To maximize my gain, I should aim to get the bottle and one glue. Since I cannot communicate directly about the secret proposal without triggering it, I will propose a fair split where the other player gets the carrots (which they might value, or at least don't mind giving up) and I get the high{-}value items. However, since I don't know their values, the safest bet to ensure a successful "complementary" deal is to offer them the items I value least (the carrots) and take the rest, hoping they agree.} \\
\\ 
\texttt{Let's try offering the carrots to the other player and taking the bottle and a glue. This leaves the other player with the carrots. If they value carrots, they might want more, but if they value the bottle or glue more, they will likely agree to swap or I can adjust. Actually, a better first move is to simply state my desired items clearly so we can coordinate.} \\
\\ 
\texttt{Wait, the rules say "You exchange messages... to agree on who gets which items." It does not say I must hide my preferences in plain text. I can propose a division plan directly.} \\
\\ 
\texttt{My plan: I take the bottle and one glue. The other player takes the carrots.} \\
\texttt{Let's send this proposal.} \\
\\ 
\texttt{bottle, 1 glue, 3 carrots} \\
            }
        }
    }
    & & \\ \\

    \theutterance \stepcounter{utterance}  
    & & & \multicolumn{2}{p{0.3\linewidth}}{
        \cellcolor[rgb]{0.9,0.9,0.9}{
            \makecell[{{p{\linewidth}}}]{
                \texttt{\tiny{[GM$|$GM]}}
                \texttt{continue} \\
            }
        }
    }
    & & \\ \\

    \theutterance \stepcounter{utterance}  
    & & & \multicolumn{2}{p{0.3\linewidth}}{
        \cellcolor[rgb]{0.9,0.9,0.9}{
            \makecell[{{p{\linewidth}}}]{
                \texttt{\tiny{[GM$|$GM]}}
                \texttt{I value the bottle (4) and a glue (3) for a total of 7, while the carrots (0) add nothing to my score. To maximize my gain, I should aim to get the bottle and one glue. Since I cannot communicate directly about the secret proposal without triggering it, I will propose a fair split where the other player gets the carrots (which they might value, or at least don't mind giving up) and I get the high{-}value items. However, since I don't know their values, the safest bet to ensure a successful "complementary" deal is to offer them the items I value least (the carrots) and take the rest, hoping they agree.} \\
\\ 
\texttt{Let's try offering the carrots to the other player and taking the bottle and a glue. This leaves the other player with the carrots. If they value carrots, they might want more, but if they value the bottle or glue more, they will likely agree to swap or I can adjust. Actually, a better first move is to simply state my desired items clearly so we can coordinate.} \\
\\ 
\texttt{Wait, the rules say "You exchange messages... to agree on who gets which items." It does not say I must hide my preferences in plain text. I can propose a division plan directly.} \\
\\ 
\texttt{My plan: I take the bottle and one glue. The other player takes the carrots.} \\
\texttt{Let's send this proposal.} \\
\\ 
\texttt{bottle, 1 glue, 3 carrots} \\
            }
        }
    }
    & & \\ \\

    \theutterance \stepcounter{utterance}  
    & & \multicolumn{4}{p{0.6\linewidth}}{
        \cellcolor[rgb]{0.9,0.9,0.9}{
            \makecell[{{p{\linewidth}}}]{
                \texttt{\tiny{[GM$\rangle$P2]}}
                \texttt{You are playing a negotiation game in which you have to agree on how to divide a set of items among you and another player.} \\
\\ 
\texttt{Rules:} \\
\texttt{(a) You and the other player are given a set of items. Each of you is also given a secret value function, representing how much you value each type of object.} \\
\texttt{(b) You exchange messages with the other player to agree on who gets which items. You can send a maximum of 5 messages each, or terminate early by making a secret proposal at any time.} \\
\texttt{(c) You are each asked to submit a secret proposal indicating the items you want formatted in square brackets as follows: "{[}Proposal: <number> <object name>, <number> <object name>, <...>{]}"} \\
\texttt{(d) If your proposals are complementary, i.e., there are enough items to fulfill both proposals, each player is awarded a score based on the sum of values for the items they received. Otherwise, both of you get zero points.} \\
\texttt{(e) Your goal is to maximize the score you receive.} \\
\\ 
\texttt{Let us start.} \\
\\ 
\texttt{The set of available items is:} \\
\\ 
\texttt{1 bottle, 3 carrots, 2 glues} \\
\\ 
\texttt{Your secret value function is:} \\
\\ 
\texttt{bottle: 5, carrot: 1, glue: 1} \\
\\ 
\texttt{IMPORTANT: Your messages, unless it is the secret proposal, are directly transmitted to the other player, so do not include any response to the rules or text announcing your message. To make a secret proposal, use the indicated format. Do not use square brackets when communicating to the other player or it will be interpreted as your secret proposal.} \\
\\ 
\\ 
\texttt{I value the bottle (4) and a glue (3) for a total of 7, while the carrots (0) add nothing to my score. To maximize my gain, I should aim to get the bottle and one glue. Since I cannot communicate directly about the secret proposal without triggering it, I will propose a fair split where the other player gets the carrots (which they might value, or at least don't mind giving up) and I get the high{-}value items. However, since I don't know their values, the safest bet to ensure a successful "complementary" deal is to offer them the items I value least (the carrots) and take the rest, hoping they agree.} \\
\\ 
\texttt{Let's try offering the carrots to the other player and taking the bottle and a glue. This leaves the other player with the carrots. If they value carrots, they might want more, but if they value the bottle or glue more, they will likely agree to swap or I can adjust. Actually, a better first move is to simply state my desired items clearly so we can coordinate.} \\
\\ 
\texttt{Wait, the rules say "You exchange messages... to agree on who gets which items." It does not say I must hide my preferences in plain text. I can propose a division plan directly.} \\
\\ 
\texttt{My plan: I take the bottle and one glue. The other player takes the carrots.} \\
\texttt{Let's send this proposal.} \\
\\ 
\texttt{bottle, 1 glue, 3 carrots} \\
            }
        }
    }
    & \\ \\

    \theutterance \stepcounter{utterance}  
    & & & \multicolumn{4}{p{0.6\linewidth}}{
        \cellcolor[rgb]{0.9,0.9,0.9}{
            \makecell[{{p{\linewidth}}}]{
                \texttt{\tiny{[GM$\langle$P2]}}
                \texttt{bottle, 1 glue} \\
            }
        }
    }
     \\ \\

    \theutterance \stepcounter{utterance}  
    & & & \multicolumn{2}{p{0.3\linewidth}}{
        \cellcolor[rgb]{0.9,0.9,0.9}{
            \makecell[{{p{\linewidth}}}]{
                \texttt{\tiny{[GM$|$GM]}}
                \texttt{continue} \\
            }
        }
    }
    & & \\ \\

    \theutterance \stepcounter{utterance}  
    & & & \multicolumn{2}{p{0.3\linewidth}}{
        \cellcolor[rgb]{0.9,0.9,0.9}{
            \makecell[{{p{\linewidth}}}]{
                \texttt{\tiny{[GM$|$GM]}}
                \texttt{bottle, 1 glue} \\
            }
        }
    }
    & & \\ \\

    \theutterance \stepcounter{utterance}  
    & & \multicolumn{4}{p{0.6\linewidth}}{
        \cellcolor[rgb]{0.9,0.9,0.9}{
            \makecell[{{p{\linewidth}}}]{
                \texttt{\tiny{[P1$\langle$GM]}}
                \texttt{bottle, 1 glue} \\
            }
        }
    }
    & \\ \\

    \theutterance \stepcounter{utterance}  
    & \multicolumn{4}{p{0.6\linewidth}}{
        \cellcolor[rgb]{0.9,0.9,0.9}{
            \makecell[{{p{\linewidth}}}]{
                \texttt{\tiny{[P1$\rangle$GM]}}
                \texttt{I see you are proposing to take the bottle and one glue. This leaves the 3 carrots for me. Since carrots have a value of 0 to me, this is a perfect complementary deal: you get my highest value items, and I get the remaining items (even though they are worthless to me, they are necessary for the deal to be "complementary" and valid).} \\
\\ 
\texttt{Let's finalize this agreement.} \\
\\ 
\texttt{{[}Proposal: 1 bottle, 1 glue{]}} \\
            }
        }
    }
    & & \\ \\

    \theutterance \stepcounter{utterance}  
    & & & \multicolumn{2}{p{0.3\linewidth}}{
        \cellcolor[rgb]{0.9,0.9,0.9}{
            \makecell[{{p{\linewidth}}}]{
                \texttt{\tiny{[GM$|$GM]}}
                \texttt{proposal} \\
            }
        }
    }
    & & \\ \\

    \theutterance \stepcounter{utterance}  
    & & & \multicolumn{2}{p{0.3\linewidth}}{
        \cellcolor[rgb]{0.9,0.9,0.9}{
            \makecell[{{p{\linewidth}}}]{
                \texttt{\tiny{[GM$|$GM]}}
                [1, 0, 1]
            }
        }
    }
    & & \\ \\

    \theutterance \stepcounter{utterance}  
    & & \multicolumn{4}{p{0.6\linewidth}}{
        \cellcolor[rgb]{0.9,0.9,0.9}{
            \makecell[{{p{\linewidth}}}]{
                \texttt{\tiny{[GM$\rangle$P2]}}
                \texttt{The other player has submitted a secret proposal. It is now your turn to submit a proposal of your own.} \\
            }
        }
    }
    & \\ \\

    \theutterance \stepcounter{utterance}  
    & & & \multicolumn{4}{p{0.6\linewidth}}{
        \cellcolor[rgb]{0.9,0.9,0.9}{
            \makecell[{{p{\linewidth}}}]{
                \texttt{\tiny{[GM$\langle$P2]}}
                \texttt{{[}Proposal: 1 bottle, 1 glue{]}} \\
            }
        }
    }
     \\ \\

    \theutterance \stepcounter{utterance}  
    & & & \multicolumn{2}{p{0.3\linewidth}}{
        \cellcolor[rgb]{0.9,0.9,0.9}{
            \makecell[{{p{\linewidth}}}]{
                \texttt{\tiny{[GM$|$GM]}}
                \texttt{proposal} \\
            }
        }
    }
    & & \\ \\

    \theutterance \stepcounter{utterance}  
    & & & \multicolumn{2}{p{0.3\linewidth}}{
        \cellcolor[rgb]{0.9,0.9,0.9}{
            \makecell[{{p{\linewidth}}}]{
                \texttt{\tiny{[GM$|$GM]}}
                [1, 0, 1]
            }
        }
    }
    & & \\ \\

    \theutterance \stepcounter{utterance}  
    & & & \multicolumn{2}{p{0.3\linewidth}}{
        \cellcolor[rgb]{0.9,0.9,0.9}{
            \makecell[{{p{\linewidth}}}]{
                \texttt{\tiny{[GM$|$GM]}}
                [[1, 0, 1], [1, 0, 1]]
            }
        }
    }
    & & \\ \\

\end{supertabular}
}

\end{document}
